\subsection{Cloudové API}
\label{subsec:deployment-cloud}

Cloudové API představuje způsob přístupu k jazykovému modelu provozovanému na serverech externího poskytovatele.
Systém Kelvin komunikuje s modelem prostřednictvím HTTP požadavků, tedy odešle vstupní data (prompt, zadání a zdrojový kód) na API endpoint poskytovatele a obdrží vygenerovanou odpověď.
Veškerý výpočet probíhá na straně poskytovatele, takže pro provoz není potřeba žádný specializovaný hardware na školní infrastruktuře.

Mezi hlavní poskytovatele cloudových API pro velké jazykové modely patří:
\begin{itemize}
    \item \textbf{OpenAI} -- modely řady GPT-5.4, GPT-5-mini a další varianty\cite{openai_api}.
    \item \textbf{Anthropic} -- modely Claude Opus~4.6, Claude Sonnet~4.6 a Claude Haiku~4.5\cite{anthropic_api}.
    \item \textbf{Google} -- modely Gemini dostupné prostřednictvím Google AI Studio nebo Vertex AI~\cite{google_vertex}.
    \item \textbf{Mistral AI} -- modely Mistral Large a Mistral Nemo, přičemž některé jsou k dispozici i jako open-source ke stažení~\cite{mistral_api}.
\end{itemize}

\subsubsection*{Hardwarové nároky}
\label{subsubsec:deployment-cloud-hardware}

Z pohledu systému Kelvin jsou hardwarové nároky pro použití cloudového API prakticky nulové.
Na straně Kelvinu je potřeba pouze stabilní internetové připojení a server schopný zpracovávat HTTP požadavky.
Školní infrastruktura, která leží na optické síti CETIN, tuto podmínku bez problémů splňuje.

\subsubsection*{Limity}
\label{subsubsec:deployment-cloud-limits}

Cloudové API zavádějí několik typů omezení, které je nutné při návrhu systému zohlednit:

\begin{enumerate}
    \item \textbf{Limit na počet tokenů v kontextu} -- každý model má maximální délku zpracovatelného vstupu i výstupu.
    U moderních modelů činí tento limit zpravidla 128~000 tokenů nebo více, což je pro potřeby Kelvinu dostatečné.
    Systém však musí zajistit, aby požadavky nepřekročily tento limit, a případnou chybovou odpověď API korektně ošetřit.

    \item \textbf{Limit na počet požadavků} -- poskytovatelé zavádějí takzvané \textit{rate limits}, které omezují počet požadavků v daném časovém období (například 100 požadavků za minutu).
    V kontextu Kelvinu je toto omezení relevantní zejména v období před termínem odevzdání, kdy může počet současných požadavků výrazně vzrůst.
    Systém proto musí implementovat mechanismus řízení zátěže.

    \item \textbf{Finanční náklady} -- cloudové API je zpoplatněno na základě počtu zpracovaných tokenů, což při rozsáhlém nasazení představuje opakující se výdaj.
    Podrobná analýza nákladů pro jednotlivé modely je uvedena níže.

    \item \textbf{Závislost na třetí straně} -- systém Kelvin se stává závislým na dostupnosti a spolehlivosti poskytovatele.
    Výpadek služby, změna API rozhraní nebo ukončení podpory konkrétního modelu mohou narušit funkčnost celého systému.
\end{enumerate}

\subsubsection*{Ochrana dat}
\label{subsubsec:deployment-cloud-data-protection}

Otázka ochrany dat studentů při použití cloudového API byla podrobně rozebrána v~\secref[Sekci]{subsubsec:gdpr-data-transfer}.
V případě cloudového nasazení studentské zdrojové kódy a zadání úloh opouštějí školní infrastrukturu a jsou zpracovávány na serverech poskytovatele.
Tato skutečnost může být rozhodujícím faktorem pro volbu mezi cloudovým API a on-premise nasazením.

\subsubsection*{Náklady}
\label{subsubsec:deployment-cloud-costs}

Náklady cloudových API jsou určeny cenovým modelem \textit{za token}.
Poskytovatelé zpravidla rozlišují vstupní tokeny (text odesílaný modelu) a výstupní tokeny (text generovaný modelem), přičemž výstupní tokeny bývají výrazně dražší.

% ===================================================================
%   Chain-of-thought přístup (3 kroky: Draft analysis, Critique analysis, Review analysis):
%     Průměr vstupních tokenů:  9 053  (min: 7 247, max: 9 687)
%     Průměr výstupních tokenů: 5 149  (min: 3 863, max: 6 404)
%
%   One-shot přístup (1 krok: One-shot analysis):
%     Průměr vstupních tokenů:  2 043  (min:   891, max: 4 508)
%     Průměr výstupních tokenů:   501  (min:   305, max:   590)
% ===================================================================
\newcommand{\chainInToks}{9053}
\newcommand{\chainOutToks}{5149}
\newcommand{\oneshotInToks}{2043}
\newcommand{\oneshotOutToks}{501}

% Výpočet nákladů: #1 = cena vstupu ($/1M tokenů), #2 = cena výstupu ($/1M tokenů)
%                  #3 = počet vstupních tokenů,    #4 = počet výstupních tokenů
\newcommand{\computecost}[4]{%
  \begingroup
  \pgfkeys{/pgf/fpu, /pgf/fpu/output format=fixed}%
  \pgfmathparse{(#3 * (#1 / 1000000)) + (#4 * (#2 / 1000000))}%
  \pgfmathsetmacro{\myresult}{\pgfmathresult}%
  \pgfkeys{/pgf/fpu=false}%
  \pgfkeys{/pgf/number format/.cd, use comma, fixed, precision=4}%
  \pgfmathprintnumber{\myresult}%
  \endgroup
}

% Výpočet nákladů pro Chain-of-Thought
% #1 = cena vstupu ($/1M), #2 = cena výstupu ($/1M)
\newcommand{\computeCostCOT}[2]{%
  \computecost{#1}{#2}{\chainInToks}{\chainOutToks}%
}

% Výpočet nákladů pro One-shot
% #1 = cena vstupu ($/1M), #2 = cena výstupu ($/1M)
\newcommand{\computeCostOS}[2]{%
  \computecost{#1}{#2}{\oneshotInToks}{\oneshotOutToks}%
}

\begin{table}[H]
    \centering
    \resizebox{\textwidth}{!}{%
        \begin{tabular}{lllll}
            \toprule
            & & & \multicolumn{2}{c}{\textbf{Náklady na odevzdání}} \\
            \textbf{Model} & \textbf{Vstup (\$/1M)} & \textbf{Výstup (\$/1M)} & \textbf{chain-of-thought (\$)} & \textbf{one-shot (\$)} \\
            \midrule
                GPT-5.4\cite{gpt-5.4}                               & 2.50  & 15.00  & $\sim\computeCostCOT{2.50}{16.00}$ & $\sim\computeCostOS{2.50}{15.00}$ \\
                GPT-5.4-mini\cite{gpt-5.4-mini}                     & 0.75  & 4.50   & $\sim\computeCostCOT{0.75}{4.50}$  & $\sim\computeCostOS{0.75}{4.50}$ \\
                GPT-5.4-nano\cite{gpt-5.4-nano}                     & 0.20  & 1.25   & $\sim\computeCostCOT{0.20}{1.25}$  & $\sim\computeCostOS{0.20}{1.25}$ \\
            \midrule
                Claude 4.6 Opus\cite{claude-opus-4.6}               & 5.00  & 25.00  & $\sim\computeCostCOT{5.00}{25.00}$ & $\sim\computeCostOS{5.00}{25.00}$ \\
                Claude 4.6 Sonnet\cite{claude-sonnet-4.6}           & 3.00  & 15.00  & $\sim\computeCostCOT{3.00}{15.00}$ & $\sim\computeCostOS{3.00}{15.00}$ \\
                Claude 3.5 Haiku\cite{claude-3.5}                   & 0.80  & 4.00   & $\sim\computeCostCOT{0.80}{4.00}$  & $\sim\computeCostOS{0.80}{4.00}$ \\
                Claude 3.0 Haiku\cite{claude-3}                     & 0.25  & 1.25   & $\sim\computeCostCOT{0.25}{1.25}$  & $\sim\computeCostOS{0.25}{1.25}$ \\
            \midrule
                Gemini 3.1 Pro\cite{geminy-3.1-pro}                 & 2.00  & 12.00  & $\sim\computeCostCOT{2.00}{12.00}$ & $\sim\computeCostOS{2.00}{12.00}$ \\
                Gemini 3.1 Flash-Lite\cite{geminy-3.1-flash-lite}   & 0.25  & 1.50   & $\sim\computeCostCOT{0.25}{1.50}$  & $\sim\computeCostOS{0.25}{1.50}$ \\
                Gemini 3 Flash\cite{geminy-3-flash}                 & 0.50  & 3.00   & $\sim\computeCostCOT{0.50}{3.00}$  & $\sim\computeCostOS{0.50}{3.00}$ \\
                Gemini 2.5 Flash\cite{geminy-2.5-flash}             & 0.30  & 2.50   & $\sim\computeCostCOT{0.30}{2.50}$  & $\sim\computeCostOS{0.30}{2.50}$ \\
            \bottomrule
        \end{tabular}
    }
    \caption{Přibližné ceny vybraných cloudových API modelů. Náklady jsou vypočítány pro dva přístupy: chain-of-thought (CoT) a one-shot (OS), na základě průměrného počtu vstupních a výstupních tokenů pro jedno odevzdání.}
    \label{tab:cloud-api-costs}
\end{table}

Z pohledu nákladů se jako nejvýhodnější jeví modely GPT-5.4-nano a Claude 3.0 Haiku, které vychází na méně než \$0,002 za odevzdání při one-shot přístupu.
Použití chain-of-thought přístupu se třemi kroky analýzy je výrazně nákladnější: modely střední třídy (Claude 3.5 Haiku, GPT-5.4-mini) vychází na přibližně \$0,028--\$0,030 za odevzdání, oproti méně než \$0,004 při one-shot přístupu.
Výjimkou je Gemini 3.1 Flash-Lite, který i při chain-of-thought přístupu dosahuje nákladů okolo \$0,01 za odevzdání.

Nejvýkonnější modely (Claude 4.6 Opus, Claude 4.6 Sonnet, GPT-5.4) přesahují \$0,10 za odevzdání u CoT přístupu, přičemž Claude 4.6 Opus dosahuje nejvyšších nákladů \$0,174, což je pro rozsáhlé školní nasazení výrazně nákladnější.
Pro konkrétní volbu modelu bude rozhodující výsledek srovnávacího testování kvality výstupů.
Z pohledu nákladů vychází nejlépe menší modely, pokud jejich analytická kvalita pro dané požadavky je dostatečná.

\endinput